\clearpage
\section{Organisation of the Stacked L3 Arrays}
\label{app:l3d-organisation}

This appendix collects the arithmetic on the organisation of the 64\,MB stacked L3 die (L3D) summarised in \cref{sec:cache-interface}.

\textbf{How the stacked arrays are organised.} The L3D contains 512 data macros of 128\,KB each and 1088 tag and LRU macros of 6\,KB each \cite[slide~25]{wuu2022vcache}. The first figure checks: $512 \times 128\,\text{KB} = 64\,\text{MB}$. The arrays use a high-density SRAM bit cell and two supply rails, one for the bit cells and one for the peripheral circuits, with ``extensive power reduction features'' \cite[slide~25]{wuu2022vcache}. The L3 is protected by DECTED (double-error-correcting, triple-error-detecting) error-correcting code (ECC) \cite[slide~19]{wuu2022vcache}.

\textbf{The cost of tags.} The tag and LRU macros hold $1088 \times 6\,\text{KB} = 6528$\,KB, about 10\% of the 64\,MB of data. With the 64-byte lines used throughout Zen~3 \cite{amd2022epyc7003}, the L3D holds $64\,\text{MB}/64\,\text{B} = 1\,048\,576$ lines, so the tag, state and replacement storage averages $6528 \times 1024 \times 8 / 1\,048\,576 \approx 51$ bits per line. The address split of \cref{sec:cache-org} shows where those bits go. A 64\,MB, 16-way cache with 64-byte lines has $2^{26}/(2^{6} \times 16) = 65\,536 = 2^{16}$ sets, so 6 address bits select the byte and 16 select the set. With the 48-bit physical addresses AMD documents for Zen~3 EPYC processors \cite[p.~10]{advancedmicrodevicesinc.UsingSEVAMD2023}, each tag would need $48 - 16 - 6 = 26$ bits, leaving roughly 25 bits per line for coherence state, replacement information and error-correction bits. The overhead is significant: a cache die is not all data storage.

\textbf{An open detail.} The combined 96\,MB L3 is 16-way set associative \cite[slide~20]{wuu2022vcache}. With 64-byte lines that is $96\,\text{MB} / (64\,\text{B} \times 16) = 98\,304 = 3 \times 2^{15}$ sets, which is not a power of two, so the simple indexing of \cref{sec:cache-org}, in which a fixed group of address bits selects the set, cannot cover it directly. Because each die's portion of the L3 has its own data, tag and LRU arrays \cite[slide~20]{wuu2022vcache}, the base and stacked portions may be indexed separately, but AMD's materials do not say how addresses are divided between them. Lam, examining the Zen~4 version, describes each nominally 4\,MB L3 slice as having its own interface to an 8\,MB extension on the cache die, with its own tags and LRU state, and suggests that the controller \emph{probably} chooses between the base and stacked portions using a subset of the address bits \cite{lamAMDs7950X3DZen2023}. That arrangement would split each slice's lines between the two dies in a 1:2 ratio and would explain the non-power-of-two total, but it remains an informed inference rather than an AMD statement.
